%% ma: L12-B11-0020
%% lop: 12
%% chuong: 4
%% bai: 11
%% dang: 12.11.15
%% loai: DS
%% muc: [TH, TH, VD, VD]
%% dap_an: "SSĐĐ"
%% nguon: "Ảnh giáo viên gửi 24/10/2026 (M06): Kỳ thi thử đợt IV, Thi thử TN THPT 2027 (Biên soạn: Aimbet), Phần 2 Câu 1"
%% kien_thuc: "Bài 5 (ứng dụng đạo hàm: lợi nhuận, tốc độ thay đổi), Bài 2 (giá trị lớn nhất), Bài 11 (tìm E(x) từ E'(x) bằng nguyên hàm)"
%% ghi_chu: "Xếp ở Bài 11 vì phải tìm E(x) bằng nguyên hàm (Bài 11 sau Bài 5); đáp án tự giải, đề không kèm đáp án"
%% trang_thai: chinh-thuc
\begin{ex}
Trong một phân xưởng của một nhà máy hóa chất, người ta sản xuất một loại hóa chất để cung cấp cho các doanh nghiệp. Chi phí sản xuất $x$ kg hóa chất đó $C(x)$ (đơn vị: nghìn đồng) được xác định bởi hàm số $C(x)=0{,}1x^3-7{,}5x^2+170x+2000$. Sau khi khảo sát thị trường, giá bán của mỗi kg hóa chất đó là $350$ nghìn đồng/kg. Khi sản lượng của hóa chất đó vượt quá $80$ kg, phân xưởng phải trả thêm thuế môi trường $E(x)$ có tốc độ thay đổi (đơn vị: nghìn đồng/kg) được xác định bởi hàm số $E'(x)=2x-150$ và $E(80)=0$. Lợi nhuận $P(x)$ được xác định bởi công thức $P(x)=R(x)-C(x)-E(x)$, với $R(x)$ là doanh thu.
\choiceTF
{Khi sản lượng của hóa chất đó dưới $80$ kg, tốc độ thay đổi của lợi nhuận $P'(x)$ được xác định bởi hàm số $P'(x)=-0{,}3x^2+18x+150$.}
{Tổng lợi nhuận của phân xưởng luôn luôn tăng trên đoạn $[50;70]$.}
{\True Tổng lợi nhuận của phân xưởng đạt giá trị lớn nhất bằng $14$ triệu $200$ nghìn đồng.}
{\True Khi phân xưởng sản xuất $90$ kg hóa chất đó, tổng lợi nhuận nhỏ hơn $2$ triệu đồng.}
\loigiai{Với $0\le x\le80$: $P(x)=350x-C(x)$, $P'(x)=-0{,}3x^2+15x+180$, nên a) sai.
$P'(x)=0\Leftrightarrow x=60$ (loại $x=-10$); $P'>0$ trên $[50;60)$, $P'<0$ trên $(60;70]$, nên b) sai.
Với $x>80$: $E(x)=x^2-150x+5600$, $P'(x)=-0{,}3x^2+13x+330<0$ (vì $P'(80)=-550$ và $P'$ giảm khi $x>80$), nên $P$ giảm trên $[80;+\infty)$.
Vậy $\max P=P(60)=21\,000-6\,800=14\,200$ nghìn đồng: c) đúng.
$P(90)=31\,500-29\,450-200=1\,850$ nghìn đồng $<2$ triệu đồng: d) đúng.}
\end{ex}
